\chapter{System Overview}
\label{ch:system-overview}

\textit{\sysname{} is a full-stack system that spans three tiers, five
subsystems, and three databases.  This chapter maps the terrain so that the
deep dives in Part~III have a frame of reference.}

\section{Architecture}
\label{sec:architecture}

At the highest level, \sysname{} follows a three-tier architecture: a
React-based frontend, a Python API server, and a persistence layer composed of
three specialized databases plus a caching tier.
Figure~\ref{fig:architecture} shows the major components and the protocols that
connect them.

\begin{figure}[ht]
\centering
\begin{tikzpicture}[
    tier/.style={draw, rounded corners, minimum width=12cm, minimum height=1.6cm,
                 fill=ContinuumViolet!8, font=\small},
    dbbox/.style={draw, rounded corners, minimum width=3.4cm, minimum height=1.4cm,
                  fill=ContinuumCyan!10, font=\small, align=center},
    arr/.style={-{Stealth[length=5pt]}, thick, ContinuumGray},
    lbl/.style={font=\scriptsize\itshape, ContinuumGray},
]

% Top tier
\node[tier] (fe) at (0, 4)
    {\textbf{Next.js 16} (React 19, App Router, shadcn/ui) --- Port 3000};

% Middle tier
\node[tier] (be) at (0, 2)
    {\textbf{FastAPI} (Python 3.12+, async, Pydantic v2) --- Port 8000};

% Bottom tier databases
\node[dbbox] (pg) at (-4.3, 0)
    {\textbf{PostgreSQL}\\ Port 5432};
\node[dbbox] (neo) at (0, 0)
    {\textbf{Neo4j}\\ Ports 7474 / 7687};
\node[dbbox] (redis) at (4.3, 0)
    {\textbf{Redis}\\ Port 6379};

% Arrows
\draw[arr] (fe.south) -- (be.north)
    node[lbl, midway, right=3pt] {REST / SSE / WebSocket};
\draw[arr] (be.south) -- (pg.north)
    node[lbl, midway, left=2pt] {SQLAlchemy async};
\draw[arr] (be.south) -- (neo.north)
    node[lbl, midway, right=2pt] {Bolt / Cypher};
\draw[arr] (be.south) -- (redis.north)
    node[lbl, midway, right=2pt] {aioredis};

% Tier labels
\node[font=\scriptsize\bfseries, ContinuumViolet, anchor=east] at (-6.5, 4) {FRONTEND};
\node[font=\scriptsize\bfseries, ContinuumViolet, anchor=east] at (-6.5, 2) {BACKEND};
\node[font=\scriptsize\bfseries, ContinuumViolet, anchor=east] at (-6.5, 0) {DATA};

\end{tikzpicture}
\caption{Three-tier architecture of \sysname{}.}
\label{fig:architecture}
\end{figure}

The frontend communicates with the backend over three protocols: standard REST
for CRUD operations, Server-Sent Events (SSE) for streaming AI responses in
the Ask interface, and WebSocket for real-time conversation capture.  The
backend talks to three databases, each chosen for a specific strength discussed
in Section~\ref{sec:tech-stack}.


\section{Five Subsystems}
\label{sec:five-subsystems}

\sysname{} is organized into five subsystems, each responsible for one stage of
the pipeline from raw conversation to searchable knowledge.
Table~\ref{tab:subsystems} summarizes them.

\begin{table}[ht]
\centering
\caption{The five subsystems of \sysname{}.}
\label{tab:subsystems}
\small
\begin{tabularx}{\textwidth}{l X l l}
\toprule
\textbf{Subsystem} & \textbf{Purpose} & \textbf{Key Files} & \textbf{Technologies} \\
\midrule
Capture &
    Passive log extraction, AI interviews &
    \codefile{capture.py} &
    WebSocket, watcher \\
Extraction &
    LLM decision trace extraction &
    \codefile{extractor.py} &
    NVIDIA NIM \\
Resolution &
    7-stage entity deduplication &
    \codefile{entity\_resolver.py} &
    RapidFuzz, Redis \\
Retrieval &
    Hybrid search + GraphRAG &
    \codefile{graph\_rag.py} &
    Neo4j, RRF \\
Integration &
    AI agent access via MCP &
    \codefile{agent.py} &
    MCP (5 tools) \\
\bottomrule
\end{tabularx}
\normalsize
\end{table}

\textbf{Capture} watches for new Claude Code conversation logs on disk and
accepts real-time input via WebSocket.  It does not interpret the content---it
simply ensures that raw conversations reach the extraction stage.

\textbf{Extraction} sends each conversation to a 49-billion-parameter language
model (Llama 3.3 Nemotron Super via NVIDIA NIM) with a structured prompt that
requests decision traces in JSON format.  The model identifies triggers,
context, options, decisions, and rationale.

\textbf{Resolution} takes the raw entity mentions from extraction (e.g.,
``postgres'', ``Postgres'', ``PostgreSQL'') and resolves them to canonical
graph nodes through a seven-stage pipeline that progresses from cheap checks
(cache, exact match, canonical lookup) to expensive ones (fuzzy matching,
embedding similarity).

\textbf{Retrieval} provides two search modes: keyword-based hybrid search
(combining fulltext and vector similarity with Reciprocal Rank Fusion) and
conversational GraphRAG that expands seed results into subgraphs for
context-rich answers.

\textbf{Integration} exposes the knowledge graph to external AI agents via the
Model Context Protocol (MCP), providing five tools that let agents search
decisions, query entities, and explore relationships.


\section{Data Flow}
\label{sec:data-flow}

Figure~\ref{fig:data-flow} traces a single conversation from capture to
queryable knowledge.

\begin{figure}[ht]
\centering
\begin{tikzpicture}[
    box/.style={draw, rounded corners, minimum width=2.8cm, minimum height=1cm,
                fill=ContinuumViolet!8, font=\small, align=center},
    db/.style={draw, rounded corners, minimum width=2.2cm, minimum height=1cm,
               fill=ContinuumCyan!10, font=\small, align=center},
    arr/.style={-{Stealth[length=5pt]}, thick, ContinuumGray},
    lbl/.style={font=\scriptsize\itshape, ContinuumGray, align=center},
]

% Nodes
\node[box] (conv) at (0, 0) {Conversation\\(JSONL log)};
\node[box] (ext) at (3.5, 0) {\codefile{extractor.py}};
\node[box] (traces) at (7, 0) {Decision\\Traces};
\node[box] (er) at (10.5, 0) {\codefile{entity\_resolver.py}};
\node[db] (neo) at (10.5, -2) {\textbf{Neo4j}};
\node[box] (rag) at (7, -2) {\codefile{graph\_rag.py}};
\node[box] (out) at (3.5, -2) {Search / Ask\\Interface};

% Arrows
\draw[arr] (conv) -- (ext) node[lbl, midway, above] {raw text};
\draw[arr] (ext) -- (traces) node[lbl, midway, above] {JSON};
\draw[arr] (traces) -- (er) node[lbl, midway, above] {mentions};
\draw[arr] (er) -- (neo) node[lbl, midway, right] {nodes +\\edges};
\draw[arr] (neo) -- (rag) node[lbl, midway, above] {subgraphs};
\draw[arr] (rag) -- (out) node[lbl, midway, above] {context +\\answers};

\end{tikzpicture}
\caption{Data flow from raw conversation to queryable knowledge.}
\label{fig:data-flow}
\end{figure}

\begin{enumerate}
    \item A \textbf{conversation} (stored as a JSONL log file) enters the
          system via file watcher or WebSocket.
    \item The \textbf{extractor} sends the conversation text to the LLM, which
          returns structured decision traces as JSON.
    \item Each trace contains raw entity \textbf{mentions} (e.g., ``postgres'',
          ``FastAPI'', ``Redis'').
    \item The \textbf{entity resolver} maps each mention to a canonical node
          in Neo4j, creating new nodes only when no match is found.
    \item The decision node and its \ic{INVOLVES} edges are written to
          \textbf{Neo4j}.
    \item The \textbf{GraphRAG} pipeline answers queries by finding seed nodes
          via hybrid search, expanding them into subgraphs, and synthesizing
          answers.
\end{enumerate}


\section{Multi-Tenant Architecture}
\label{sec:multi-tenant}

\sysname{} enforces strict per-user isolation.  Every decision, entity, and
relationship in the knowledge graph is scoped by a \ic{user\_id} property.
When a user queries the system, all Neo4j Cypher queries include a
\ic{WHERE n.user\_id = \$user\_id} clause, ensuring that users see only their
own decisions and entities.

This scoping is enforced at the API layer---the backend extracts the
\ic{user\_id} from the JWT token and injects it into every database query.
There is no shared namespace, no cross-user leakage, and no way for one user
to access another user's knowledge graph.

\begin{warning}
Multi-tenancy is enforced at the application layer, not the database layer.
A misconfigured query that omits the \ic{user\_id} filter would expose data
across users.  Production deployment would require additional database-level
row security policies.
\end{warning}


\section{Tech Stack}
\label{sec:tech-stack}

Each technology in \sysname{} was chosen for a specific reason.
Table~\ref{tab:tech-stack} maps the major components to their justifications.

\begin{table}[ht]
\centering
\caption{Technology choices and their rationale.}
\label{tab:tech-stack}
\begin{tabularx}{\textwidth}{l l X}
\toprule
\textbf{Technology} & \textbf{Role} & \textbf{Why This Choice} \\
\midrule
PostgreSQL 18 &
    Relational data + auth &
    Alembic migrations, async SQLAlchemy, mature ecosystem \\
Neo4j 2025.01 &
    Knowledge graph &
    Native graph traversal, Cypher queries, built-in fulltext + vector indexes \\
Redis 7.4 &
    Cache + rate limiting &
    Three cache tiers (entity, LLM, query), token-bucket rate limiting \\
NVIDIA NIM &
    LLM + embeddings &
    49B-param Nemotron for extraction, NV-EmbedQA (2048-dim) for similarity \\
Next.js 16 &
    Frontend &
    React 19, App Router, Server Components, shadcn/ui component library \\
FastAPI &
    Backend API &
    Native async, automatic OpenAPI docs, Pydantic validation \\
\bottomrule
\end{tabularx}
\end{table}

\begin{keyinsight}
A single decision touches all three databases plus Redis cache.  A new
PostgreSQL decision is stored relationally for auth context, as a Neo4j node
with entity edges for graph traversal, and cached in Redis for fast
retrieval.  Understanding where data lives is essential for debugging.
\end{keyinsight}

\begin{skipahead}
If you want to run the system before understanding every component,
skip to Chapter~\ref{ch:setup-directory-tour} for setup instructions.
\end{skipahead}
